\begin{mainquestion}{csbandwidth}{Bandwidth Limitation of a Common-Source Amplifier}

\questionpoints A common-source amplifier has small-signal gain:

\[
A_v = -g_m R_D
\]

Assume the transistor now has an output capacitance \(C_{out}\) connected at the drain node. The small signal model will look like:

\begin{center}
\begin{circuitikz}

% Input terminal and Gate-Source open circuit
\draw (-3,3) node[left]{$v_{in}$} to[short, o-] (-2,3) node[above]{G};
\draw (-2,0) node[left]{S} node[ground]{};

% Explicitly label v_gs across the open circuit
\draw (-2,3) to[open, v=$v_{gs}$] (-2,0);

% Controlled current source (Drain to Source)
% Drawn from top (3) to bottom (0) so the current points down. 
% This creates the characteristic 180-degree phase shift (the negative sign in your gain).
\draw (0,3) node[above]{D} 
      to[american controlled current source, l=$g_m v_{gs}$] (0,0) node[ground]{};

% Connect Drain to the rest of the output network
\draw (0,3) -- (4,3);

% RD branch
\draw (2,3) to[R=$R_D$] (2,0) node[ground]{};

% Cout branch
\draw (4,3) to[C=$C_{out}$] (4,0) node[ground]{};

% Output node
\draw (4,3) to[short, -o] (5,3) node[right]{$v_{out}$};

\end{circuitikz}
\end{center}

%-----------------------------------

% =========================================================
\begin{partquestion}{transferfunction}
{Derive the small-signal transfer function using sinusoidal steady-state analysis.}

\begin{answer}[2.5in]
Assume sinusoidal steady-state operation, so voltages and currents are represented by phasors.

The capacitor current is

\[
i_C = j\omega C_{out} v_{out}
\]

The resistor current is

\[
i_R = \frac{v_{out}}{R_D}
\]

The transistor small-signal drain current is

\[
i_d = g_m v_{gs}
\]

Apply KCL at the drain node (currents leaving the node through the resistor and capacitor must equal the transistor current entering the node):

\[
g_m v_{gs}
=
\frac{v_{out}}{R_D}
+
j\omega C_{out} v_{out}
\]

Factor out \(v_{out}\):

\[
g_m v_{gs}
=
v_{out}
\left(
\frac{1}{R_D}
+
j\omega C_{out}
\right)
\]

Solve for the transfer function:

\[
\frac{v_{out}}{v_{gs}}
=
\frac{g_m}
{
\frac{1}{R_D}+j\omega C_{out}
}
\]

Multiplying numerator and denominator by \(R_D\):

\[
\frac{v_{out}}{v_{gs}}
=
\frac{g_m R_D}{1+j\omega R_D C_{out}}
\]

Because a common-source amplifier inverts the signal,

\[
A_v(j\omega)
=
\frac{v_{out}}{v_{in}}
=
\frac{-g_m R_D}{1+j\omega R_D C_{out}}
\]

This is a single-pole low-pass transfer function.
\end{answer}

\pagebreak
%-----------------------------------
\end{partquestion}

% =========================================================
\begin{partquestion}{cornerfrequency}
{Derive the -3dB corner frequency.

\textit{Hint: The -3dB point occurs when}
\[
20 \log_{10} \left( \frac{|A_v(\omega_c)|}{|A_v(0)|} \right) = -3.
\]}

\begin{answer}[2.5in]
The magnitude of the transfer function is

\[
|A_v(j\omega)|
=
\frac{g_m R_D}
{\sqrt{1+(\omega R_D C_{out})^2}}
\]

At low frequency (\(\omega=0\)):

\[
|A_v(0)| = g_m R_D
\]

The -3dB point occurs when

\[
|A_v(\omega_c)|
=
\frac{|A_v(0)|}{\sqrt{2}}
\]

Substitute the magnitude expression:

\[
\frac{g_m R_D}
{\sqrt{1+(\omega_c R_D C_{out})^2}}
=
\frac{g_m R_D}{\sqrt{2}}
\]

Cancel \(g_m R_D\):

\[
\sqrt{1+(\omega_c R_D C_{out})^2}
=
\sqrt{2}
\]

Square both sides:

\[
1+(\omega_c R_D C_{out})^2 = 2
\]

\[
(\omega_c R_D C_{out})^2 = 1
\]

\[
\omega_c = \frac{1}{R_D C_{out}}
\]

Thus the corner frequency is

\[
f_c = \frac{1}{2\pi R_D C_{out}}
\]
\end{answer}

%-----------------------------------
\end{partquestion}

% =========================================================
\begin{partquestion}{bandwidth}
{Suppose $R_D = 10\,k\Omega
\quad\text{and}\quad
C_{out} = 2\,pF$. Calculate the -3dB corner frequency.}

\begin{answer}[2.5in]
Given

\[
R_D = 10\,k\Omega
\]

\[
C_{out} = 2\,pF
\]

The -3dB frequency is

\[
f_c =
\frac{1}{2\pi R_D C_{out}}
\]

\[
f_c =
\frac{1}{2\pi (10^4)(2\times10^{-12})}
\]

\[
f_c =
\frac{1}{2\pi (2\times10^{-8})}
\]

\[
f_c =
\frac{1}{1.26\times10^{-7}}
\]

\[
f_c \approx 7.96\,\text{MHz}
\]

\[
\boxed{f_c \approx 8\,\text{MHz}}
\]
\end{answer}

%-----------------------------------
\end{partquestion}

% =========================================================
\begin{partquestion}{gaintradeoff}
{Explain why increasing \(R_D\) (to increase gain) reduces bandwidth.}

\begin{answer}[2in]
The low-frequency gain of the amplifier is

\[
A_v = -g_m R_D
\]

Thus increasing \(R_D\) increases the gain.

However, the pole frequency is

\[
f_c = \frac{1}{2\pi R_D C_{out}}
\]

If \(R_D\) increases, the RC time constant

\[
\tau = R_D C_{out}
\]

also increases, which lowers the corner frequency.

Therefore increasing gain by increasing \(R_D\) reduces the bandwidth.  
This is a fundamental gain–bandwidth tradeoff in amplifier design.
\end{answer}

\end{partquestion}

\end{mainquestion}
